\documentclass[11pt]{article}
\usepackage[margin=1in]{geometry}
\usepackage{amsmath}
\usepackage{hyperref}

\title{Affine Schedule Response in a Synthetic Rotated-Surface-Code Fault Model}
\author{Y.Y.N. Li}
\date{v0.6.1, 2026-08-02}

\begin{document}
\maketitle

\begin{abstract}
This note documents a prospective, fail-closed numerical mechanism experiment
for an affine schedule response in a synthetic rotated-surface-code fault
model. The experiment does not claim a threshold improvement, an asymptotic
logical-error exponent change, a compiler-derived pulse model, a hardware
advantage, or a reduction of real physical surface-code resources.
\end{abstract}

\section{Claim Boundary}
The ideal logical identity layer is inserted constructively by Stim. The
schedule does not generate distinct physical unitary evolutions. Instead, it
changes a frozen synthetic schedule-to-fault map: single-data and adjacent-pair
fault probabilities are attached to the inserted ideal layer, and PyMatching is
constructed from each arm's detector error model.

The supported statement is therefore narrow. In the frozen v0.6.0 prospective
cohort, specified and frozen before evaluation, constrained schedule
optimization reduced the minimum tested distance satisfying the frozen Wilson
criterion from \(11\) to \(9\) at target logical failure probability
\(p_L=0.025\) on all three new prospective seeds. With syndrome rounds equal
to distance, this corresponds to a 45.34\% active-coordinate qubit-round proxy
saving, not a demonstrated reduction in calibrated hardware resources.

\section{Reproducibility}
The code repository for this artifact is
\url{https://github.com/papasop/quantum}. The frozen protocol artifact is
stored as \texttt{results/v0.6.0/protocol.json}. Its SHA-256 hash under the
repository canonical JSON convention is
\[
\texttt{fec91e30001712f3d9ac84c0e45a6b70f2d5ae7189d3c9ac6d1096d47505cbf6}.
\]
The archived report certificate hash is
\[
\texttt{2db9620419ac5a7ff64510c65e0d391c4603b6c361fdd8aadd2d9f96165cbc79}.
\]
The full Monte Carlo report is regenerated by running
\texttt{python src/ft\_unit\_change\_time\_rotated\_surface\_code\_v0\_6\_0.py}
with the pinned dependencies in \texttt{requirements.txt}. The compressed
archival artifact is \texttt{results/v0.6.0/report.json.gz}.

\section{Scientific Positioning}
This is not an independent prediction validation: the optimization objective is
built from the same synthetic schedule-to-fault probability model that is later
injected into Stim. The Monte Carlo stage independently validates whether those
probabilities, after syndrome extraction and decoding, cross the declared
distance boundary. Decoded Monte Carlo outcomes are held out, but the analytic
fault map optimized by the flow is also the map injected into Stim; this is not
an out-of-model validation. The accurate positioning is a prospective,
fail-closed synthetic fault-model mechanism experiment showing that constrained
schedule optimization can cross a discrete code-distance boundary in that
model.

\end{document}
